\documentclass[a4paper,11pt]{ctexart}

% ==================== 导言区 ====================
\usepackage[left=2.5cm, right=2.5cm, top=2.5cm, bottom=2.5cm]{geometry}
\usepackage{amsmath}
\usepackage{amssymb}
\usepackage{amsfonts}
\usepackage{bm}        % 粗体数学符号
\usepackage{mathtools} % 更多数学工具
\usepackage{booktabs}  % 三线表
\usepackage{array}     % 表格列格式
\usepackage{enumitem}  % 列表定制
\usepackage{longtable} % 长表格 support

% 标题格式设置
\ctexset{
    section = {
        format = \Large\bfseries\centering,
        name = {,},
        number = \chinese{section}章,
        aftername = \hspace{0pt}
    }
}

% 自定义命令
\newcommand{\sol}{\par\noindent\textbf{解}\quad}
\newcommand{\proof}{\par\noindent\textbf{证明}\quad}
\newcommand{\ud}{\mathrm{d}}
\newcommand{\e}{\mathrm{e}}
\newcommand{\T}{\mathrm{T}}        % 转置符号 T (之前漏了这行导致报错)

\title{数值分析习题全解 (完整修正版)}
\author{}
\date{}

\begin{document}

\maketitle

% ==================== 第一章 ====================
\section*{第一章}

\begin{enumerate}[label=\textbf{1-\arabic*.}, leftmargin=*]
    \item 设 $x>0$, $x$ 的相对误差为 $\delta$, 求 $\ln x$ 的误差.
    
    \sol 设 $x$ 的近似数为 $x^*$, 则由相对误差定义知 $\left| \frac{x^*-x}{x^*} \right| = \delta$. 
    对于函数 $f(x) = \ln x$, 其导数为 $f'(x) = \frac{1}{x}$.
    根据误差传播公式, $\ln x$ 的绝对误差为:
    \[
    \varepsilon(\ln x^*) \approx |f'(x^*)| \cdot \varepsilon(x^*) = \left| \frac{1}{x^*} \right| \cdot |x^* - x| = \left| \frac{x^* - x}{x^*} \right| = \delta.
    \]
    即 $\ln x$ 的绝对误差约为 $\delta$.

    \item 设 $x$ 的相对误差为 $2\%$, 求 $x^n$ 的相对误差.
    
    \sol 设 $f(x) = x^n$. 计算函数值问题的条件数 $C_p$ 定义为:
    \[
    C_p = \left| \frac{x f'(x)}{f(x)} \right| = \left| \frac{x \cdot n x^{n-1}}{x^n} \right| = |n| = n \quad (\text{假设 } n>0).
    \]
    根据相对误差传播公式 $\varepsilon_r(f(x^*)) \approx C_p \cdot \varepsilon_r(x^*)$, 有:
    \[
    \varepsilon_r((x^*)^n) \approx n \cdot \varepsilon_r(x^*) = n \cdot 2\% = 0.02n.
    \]

    \item[5.] 计算球体积要使相对误差限为 $1\%$, 问度量半径 $R$ 时允许的相对误差限是多少?
    
    \sol 球体体积公式为 $V = \frac{4}{3}\pi R^3$. 这是一个关于 $R$ 的幂函数计算, 指数为 3.
    体积计算的条件数为:
    \[
    C_p = \left| \frac{R \cdot V'(R)}{V(R)} \right| = \left| \frac{R \cdot 4\pi R^2}{\frac{4}{3}\pi R^3} \right| = 3.
    \]
    由 $\varepsilon_r(V^*) \approx C_p \cdot \varepsilon_r(R^*)$, 得:
    \[
    \varepsilon_r(R^*) \approx \frac{1}{3} \varepsilon_r(V^*) = \frac{1}{3} \times 1\% \approx 0.33\%.
    \]
    故测量半径时允许的相对误差限约为 $0.33\%$.

    \item[9.] 正方形的边长大约为 $100 \text{ cm}$, 应怎样测量才能使其面积误差不超过 $1 \text{ cm}^2$?
    
    \sol 设正方形边长为 $x$, 面积 $A(x) = x^2$. 
    已知 $x^* \approx 100$, 要求面积的绝对误差 $|\varepsilon(A^*)| \leqslant 1$.
    由绝对误差传播公式 $\varepsilon(A^*) \approx |A'(x^*)| \cdot \varepsilon(x^*)$, 得:
    \[
    2x^* \cdot \varepsilon(x^*) \leqslant 1.
    \]
    代入 $x^* = 100$:
    \[
    200 \cdot \varepsilon(x^*) \leqslant 1 \implies \varepsilon(x^*) \leqslant \frac{1}{200} = 0.005 \text{ cm}.
    \]
    即测量边长的误差限不能超过 $0.005 \text{ cm}$.

    \item[10.] 设 $S = \frac{1}{2}gt^2$, 假定 $g$ 是准确的, 而对 $t$ 的测量有 $\pm 0.1\text{s}$ 的误差, 证明当 $t$ 增加时 $S$ 的绝对误差增加, 而相对误差却减少.
    
    \proof 
    绝对误差:
    \[
    \varepsilon(S) \approx |S'(t)| \cdot \varepsilon(t) = |gt| \cdot \varepsilon(t) = g t \cdot 0.1 = 0.1gt.
    \]
    显然, $\varepsilon(S)$ 与 $t$ 成正比, 当 $t$ 增加时, 绝对误差增加.
    
    相对误差:
    \[
    \varepsilon_r(S) = \frac{\varepsilon(S)}{S} = \frac{0.1gt}{\frac{1}{2}gt^2} = \frac{0.2}{t}.
    \]
    显然, $\varepsilon_r(S)$ 与 $t$ 成反比, 当 $t$ 增加时, 相对误差减少.
\end{enumerate}

% ==================== 第二章 ====================
\section*{第二章}

\begin{enumerate}[label=\textbf{2-\arabic*.}, leftmargin=*]
    \item 当 $x = -1, 1, 2$ 时, $f(x) = 0, -3, 4$, 求 $f(x)$ 的二次插值多项式.
    (1) 用单项式基底; (2) 用拉格朗日插值基底; (3) 用牛顿基底. 并证明三种方法得到的多项式是相同的.
    
    \sol 已知插值点 $(-1, 0), (1, -3), (2, 4)$.
    
    \textbf{(1) 单项式基底}: 设 $p_2(x) = a_0 + a_1 x + a_2 x^2$.
    代入节点得线性方程组:
    \[
    \begin{cases}
    a_0 - a_1 + a_2 = 0 \quad \dots \text{(1)} \\
    a_0 + a_1 + a_2 = -3 \quad \dots \text{(2)} \\
    a_0 + 2a_1 + 4a_2 = 4 \quad \dots \text{(3)}
    \end{cases}
    \]
    (2)-(1) 得: $2a_1 = -3 \implies a_1 = -1.5$.
    将 $a_1$ 代入 (1) 得: $a_0 + 1.5 + a_2 = 0 \implies a_0 + a_2 = -1.5$.
    将 $a_1$ 代入 (3) 得: $a_0 - 3 + 4a_2 = 4 \implies a_0 + 4a_2 = 7$.
    后两式相减: $3a_2 = 8.5 \implies a_2 = \frac{17}{6}$.
    求 $a_0$: $a_0 = -1.5 - \frac{17}{6} = -\frac{9}{6} - \frac{17}{6} = -\frac{26}{6} = -\frac{13}{3}$.
    故 \fbox{$p_2(x) = -\frac{13}{3} - \frac{3}{2}x + \frac{17}{6}x^2$}.
    
    \textbf{(2) 拉格朗日基底}:
    \begin{align*}
    l_0(x) &= \frac{(x-1)(x-2)}{(-1-1)(-1-2)} = \frac{1}{6}(x^2-3x+2) \\
    l_1(x) &= \frac{(x+1)(x-2)}{(1+1)(1-2)} = -\frac{1}{2}(x^2-x-2) \\
    l_2(x) &= \frac{(x+1)(x-1)}{(2+1)(2-1)} = \frac{1}{3}(x^2-1)
    \end{align*}
    插值多项式:
    \begin{align*}
    p_2(x) &= 0 \cdot l_0(x) + (-3) \cdot [-\frac{1}{2}(x^2-x-2)] + 4 \cdot [\frac{1}{3}(x^2-1)] \\
           &= \frac{3}{2}(x^2-x-2) + \frac{4}{3}x^2 - \frac{4}{3} \\
           &= (\frac{3}{2} + \frac{4}{3})x^2 - \frac{3}{2}x + (-3 - \frac{4}{3}) \\
           &= \frac{17}{6}x^2 - \frac{3}{2}x - \frac{13}{3}.
    \end{align*}
    结果与 (1) 相同.
    
    \textbf{(3) 牛顿基底}: 构造差商表
    \begin{center}
    \begin{tabular}{c|ccc}
    $x_i$ & $f(x_i)$ & 一阶差商 & 二阶差商 \\
    \hline
    -1 & 0 & & \\
       &   & $\frac{-3-0}{1-(-1)} = -1.5$ & \\
    1  & -3 & & $\frac{7 - (-1.5)}{2-(-1)} = \frac{8.5}{3} = \frac{17}{6}$ \\
       &   & $\frac{4-(-3)}{2-1} = 7$ & \\
    2  & 4 & & 
    \end{tabular}
    \end{center}
    牛顿插值多项式:
    \begin{align*}
    p_2(x) &= f(x_0) + f[x_0, x_1](x-x_0) + f[x_0, x_1, x_2](x-x_0)(x-x_1) \\
           &= 0 - 1.5(x+1) + \frac{17}{6}(x+1)(x-1) \\
           &= -1.5x - 1.5 + \frac{17}{6}x^2 - \frac{17}{6} \\
           &= \frac{17}{6}x^2 - \frac{3}{2}x - (\frac{9}{6} + \frac{17}{6}) \\
           &= \frac{17}{6}x^2 - \frac{3}{2}x - \frac{13}{3}.
    \end{align*}
    三种方法结果完全一致.

    \item 给出 $f(x) = \ln x$ 的数值表 (0.4, 0.5, 0.6, 0.7, 0.8), 计算 $\ln 0.54$.
    
    \sol 
    \textbf{线性插值}: 选取最近的两个点 $x_0=0.5, x_1=0.6$.
    \begin{align*}
    L_1(0.54) &= \frac{0.54-0.6}{0.5-0.6}(-0.693147) + \frac{0.54-0.5}{0.6-0.5}(-0.510826) \\
              &= 0.6 \times (-0.693147) + 0.4 \times (-0.510826) \\
              &= -0.4158882 - 0.2043304 = -0.6202186.
    \end{align*}
    
    \textbf{二次插值}: 选取 $x_0=0.4, x_1=0.5, x_2=0.6$.
    \begin{align*}
    L_2(0.54) &= \frac{(0.54-0.5)(0.54-0.6)}{(0.4-0.5)(0.4-0.6)}(-0.916291) + \frac{(0.54-0.4)(0.54-0.6)}{(0.5-0.4)(0.5-0.6)}(-0.693147) \\
              &\quad + \frac{(0.54-0.4)(0.54-0.5)}{(0.6-0.4)(0.6-0.5)}(-0.510826) \\
              &= \frac{-0.0024}{0.02}(-0.916291) + \frac{-0.0084}{-0.01}(-0.693147) + \frac{0.0056}{0.02}(-0.510826) \\
              &= -0.12(-0.916291) + 0.84(-0.693147) + 0.28(-0.510826) \\
              &= 0.10995492 - 0.58224348 - 0.14303128 \\
              &= -0.6153198.
    \end{align*}

    \item[4.] 设 $x_j$ 为互异节点, 求证: (1) $\sum x_j^k l_j(x) = x^k$; (2) $\sum (x_j-x)^k l_j(x) = 0$.
    
    \proof 
    (1) 令 $f(x) = x^k$. 当 $k \leqslant n$ 时, $f(x)$ 是次数不超过 $n$ 的多项式.
    根据插值多项式的唯一性, $n$ 次插值多项式 $L_n(x) = \sum_{j=0}^n f(x_j)l_j(x)$ 必须恒等于 $f(x)$.
    即 $\sum_{j=0}^n x_j^k l_j(x) = x^k$ 成立.
    
    (2) 令 $g(t) = (t-x)^k$, 其中 $x$ 视为固定参数, $t$ 为变量.
    $g(t)$ 是关于 $t$ 的 $k$ 次多项式. 当 $k \leqslant n$ 时, 对 $t$ 进行插值:
    $\sum_{j=0}^n g(x_j)l_j(t) = g(t)$ 对任意 $t$ 成立.
    令 $t=x$, 等式左边为 $\sum_{j=0}^n (x_j-x)^k l_j(x)$, 右边为 $g(x) = (x-x)^k = 0$ (当 $k \ge 1$). 证毕.

    \item[8.] $f(x) = x^7 + x^4 + 3x + 1$, 求 $f[2^0, \dots, 2^7]$ 及 $f[2^0, \dots, 2^8]$.
    
    \sol 差商性质: $n$ 阶差商与导数的关系为 $f[x_0, \dots, x_n] = \frac{f^{(n)}(\xi)}{n!}$.
    (1) 对于 7 次多项式 $f(x)$, 其 7 阶导数 $f^{(7)}(x) = 7!$.
    故 7 阶差商 $f[2^0, \dots, 2^7] = \frac{7!}{7!} = 1$.
    (2) 8 阶导数为 0, 故 8 阶差商 $f[2^0, \dots, 2^8] = 0$.

    \item[14.] 求次数 $\leqslant 3$ 的多项式 $P(x)$, 满足 $P(0)=0, P'(0)=1, P(1)=1, P'(1)=2$.
    
    \sol 这是一个两点三次 Hermite 插值. 节点 $x_0=0, x_1=1$.
    基函数公式:
    \begin{align*}
    \alpha_0(x) &= (1+2\frac{x-x_0}{x_1-x_0})\left(\frac{x-x_1}{x_0-x_1}\right)^2 = (1+2x)(x-1)^2 \\
    \alpha_1(x) &= (1+2\frac{x-x_1}{x_0-x_1})\left(\frac{x-x_0}{x_1-x_0}\right)^2 = (1-2(x-1))x^2 = x^2(3-2x) \\
    \beta_0(x) &= (x-x_0)\left(\frac{x-x_1}{x_0-x_1}\right)^2 = x(x-1)^2 \\
    \beta_1(x) &= (x-x_1)\left(\frac{x-x_0}{x_1-x_0}\right)^2 = (x-1)x^2
    \end{align*}
    代入数值 $f(0)=0, f(1)=1, f'(0)=1, f'(1)=2$:
    \begin{align*}
    P(x) &= 0 \cdot \alpha_0(x) + 1 \cdot \alpha_1(x) + 1 \cdot \beta_0(x) + 2 \cdot \beta_1(x) \\
         &= (3x^2-2x^3) + (x^3-2x^2+x) + (2x^3-2x^2) \\
         &= (-2+1+2)x^3 + (3-2-2)x^2 + x \\
         &= x^3 - x^2 + x.
    \end{align*}

    \item[18.] 求 $f(x) = x^2$ 在 $[a, b]$ 上的分段线性插值误差.
    
    \sol 分段线性插值 $I_h(x)$ 在子区间 $[x_i, x_{i+1}]$ 上是连接 $(x_i, x_i^2)$ 和 $(x_{i+1}, x_{i+1}^2)$ 的直线.
    误差余项公式为: $R(x) = \frac{f''(\xi)}{2}(x-x_i)(x-x_{i+1})$.
    因为 $f(x)=x^2$, 所以 $f''(x)=2$.
    误差绝对值 $|R(x)| = |(x-x_i)(x-x_{i+1})|$.
    此函数在区间中点取得最大值 $\frac{h^2}{4}$.
    故全局最大误差为 $\frac{h^2}{4}$, 其中 $h = \max(x_{i+1}-x_i)$.
\end{enumerate}

% ==================== 第三章 ====================
\section*{第三章}

\begin{enumerate}[label=\textbf{3-\arabic*.}, leftmargin=*]
    \setcounter{enumi}{3}
    \item 计算下列函数 $f(x)$ 关于 $C[0,1]$ 的范数 $\|f\|_\infty, \|f\|_1, \|f\|_2$.
    
    \sol 
    (1) $f(x) = (x-1)^3$:
    \begin{itemize}
        \item $\|f\|_\infty = \max_{x\in[0,1]} |x-1|^3 = |0-1|^3 = 1$.
        \item $\|f\|_1 = \int_0^1 |x-1|^3 \ud x = \int_0^1 (1-x)^3 \ud x = [-\frac{1}{4}(1-x)^4]_0^1 = \frac{1}{4}$.
        \item $\|f\|_2 = \sqrt{\int_0^1 (x-1)^6 \ud x} = \sqrt{[\frac{1}{7}(x-1)^7]_0^1} = \sqrt{0 - (-\frac{1}{7})} = \frac{1}{\sqrt{7}}$.
    \end{itemize}

    (2) $f(x) = |x - 1/2|$:
    \begin{itemize}
        \item $\|f\|_\infty = \max |x-0.5| = 0.5$ (在端点处取得).
        \item $\|f\|_1 = \int_0^1 |x-0.5| \ud x = 2 \int_{0.5}^1 (x-0.5) \ud x = 2 [\frac{1}{2}x^2-0.5x]_{0.5}^1 = 2(0-\frac{1}{2}(0.25-0.5)) = \frac{1}{4}$.
        \item $\|f\|_2 = \sqrt{\int_0^1 (x-0.5)^2 \ud x} = \sqrt{[\frac{1}{3}(x-0.5)^3]_0^1} = \sqrt{\frac{1}{3}(0.125 - (-0.125))} = \sqrt{\frac{0.25}{3}} = \frac{1}{\sqrt{12}}$.
    \end{itemize}

    (3) $f(x) = x^m(1-x)^n$ (简略步骤):
    \begin{itemize}
        \item $\|f\|_\infty$: 极值点在 $x = \frac{m}{m+n}$, 值为 $\frac{m^m n^n}{(m+n)^{m+n}}$.
        \item $\|f\|_1$: Beta 函数 $B(m+1, n+1) = \frac{m!n!}{(m+n+1)!}$.
        \item $\|f\|_2$: $\sqrt{B(2m+1, 2n+1)} = \sqrt{\frac{(2m)!(2n)!}{(2m+2n+1)!}}$.
    \end{itemize}

    \item[13.] 求 $f(x) = x^3$ 在 $[-1, 1]$ 上关于 $\rho(x) = 1$ 的最佳平方逼近二次多项式.
    
    \sol 勒让德多项式系 $\{P_0, P_1, P_2\}$ 为正交基.
    $P_0(x)=1, P_1(x)=x, P_2(x)=\frac{3}{2}x^2-\frac{1}{2}$.
    最佳逼近 $S_2^*(x) = c_0 P_0 + c_1 P_1 + c_2 P_2$, 系数 $c_k = \frac{(f, P_k)}{(P_k, P_k)}$.
    由于 $x^3$ 是奇函数:
    $(x^3, P_0) = \int_{-1}^1 x^3 \ud x = 0 \implies c_0=0$.
    $(x^3, P_2) = \int_{-1}^1 x^3 P_2(x) \ud x = 0 \implies c_2=0$.
    计算 $c_1$:
    $(x^3, P_1) = \int_{-1}^1 x^4 \ud x = \frac{2}{5}$.
    $(P_1, P_1) = \int_{-1}^1 x^2 \ud x = \frac{2}{3}$.
    $c_1 = \frac{2/5}{2/3} = \frac{3}{5}$.
    故最佳平方逼近多项式为 $S_2^*(x) = \frac{3}{5}x$.

    \item[14.] 求函数 $f(x)$ 在指定区间上对于 $\Phi = \text{span}\{1, x\}$ 的最佳平方逼近多项式 $y=a_0+a_1x$.
    
    \sol 通用法方程:
    \[
    \begin{pmatrix} \int_a^b 1 \ud x & \int_a^b x \ud x \\ \int_a^b x \ud x & \int_a^b x^2 \ud x \end{pmatrix} \begin{pmatrix} a_0 \\ a_1 \end{pmatrix} = \begin{pmatrix} \int_a^b f(x) \ud x \\ \int_a^b x f(x) \ud x \end{pmatrix}
    \]
    
    (1) $f(x)=1/x, [1, 3]$.
    左边矩阵: $\begin{pmatrix} 2 & 4 \\ 4 & 26/3 \end{pmatrix}$. 右边向量: $\begin{pmatrix} \ln 3 \\ 2 \end{pmatrix}$.
    解得 $a_0 \approx 1.141, a_1 \approx -0.296$.
    
    (2) $f(x)=\e^x, [0, 1]$.
    左边矩阵: $\begin{pmatrix} 1 & 0.5 \\ 0.5 & 1/3 \end{pmatrix}$. 右边向量: $\begin{pmatrix} \e-1 \\ 1 \end{pmatrix}$.
    解得 $a_0 \approx 0.873, a_1 \approx 1.690$.
    
    (3) $f(x)=\cos \pi x, [0, 1]$.
    左边矩阵同上. 右边向量: $\int_0^1 \cos \pi x = 0$, $\int_0^1 x \cos \pi x = -2/\pi^2$.
    解得 $a_0 \approx 1.216, a_1 \approx -2.431$.
    
    (4) $f(x)=\ln x, [1, 2]$.
    左边矩阵: $\begin{pmatrix} 1 & 1.5 \\ 1.5 & 7/3 \end{pmatrix}$. 右边向量: $\begin{pmatrix} 2\ln 2 - 1 \\ 2\ln 2 - 0.75 \end{pmatrix}$.
    解得 $a_0 \approx -0.637, a_1 \approx 0.682$.

    \item[17.] 用最小二乘法求 $y = a + bx^2$ 拟合数据.
    
    \sol 设 $X = x^2$, 则模型变为 $y = a + bX$.
    法方程组:
    \[
    \begin{cases}
    na + (\sum x_i^2)b = \sum y_i \\
    (\sum x_i^2)a + (\sum x_i^4)b = \sum x_i^2 y_i
    \end{cases}
    \]
    计算各求和项 ($n=5$):
    $\sum x_i^2 = 19^2+25^2+31^2+38^2+44^2 = 5327$.
    $\sum x_i^4 = 7277699$.
    $\sum y_i = 271.4$.
    $\sum x_i^2 y_i = 19^2(19)+ \dots = 369321.5$.
    代入方程组求解得: $a \approx 0.9726, b \approx 0.0500$.
    公式: $y = 0.9726 + 0.0500x^2$.
    均方误差 $\|\delta\|_2 = \sqrt{\sum (y_i - (a+bx_i^2))^2} \approx 0.130$.

    \item[18.] 化学反应数据拟合 $y = a\e^{b/t}$.
    
    \sol 线性化: $\ln y = \ln a + b(\frac{1}{t})$. 令 $Y = \ln y, X = 1/t$.
    变成线性最小二乘问题 $Y = A + bX$ (其中 $A = \ln a$).
    注意数据点 $t=0$ 处 $y=0$ 不参与对数计算, 剔除该点.
    计算得: $b \approx -7.496, A \approx 1.652 \implies a = \e^{1.652} \approx 5.215$.
    拟合公式: $y = 5.215 \e^{-7.496/t}$.
\end{enumerate}

% ==================== 第四章 ====================
\section*{第四章}

\begin{enumerate}[label=\textbf{4-\arabic*.}, leftmargin=*]
    \item 确定参数: $\int_{-h}^h f(x) \ud x \approx A_{-1} f(-h) + A_0 f(0) + A_1 f(h)$.
    
    \sol 令公式对 $f(x)=1, x, x^2$ 精确成立:
    \[
    \begin{cases}
    A_{-1} + A_0 + A_1 = \int_{-h}^h 1 \ud x = 2h \\
    -h A_{-1} + 0 + h A_1 = \int_{-h}^h x \ud x = 0 \\
    h^2 A_{-1} + 0 + h^2 A_1 = \int_{-h}^h x^2 \ud x = \frac{2}{3}h^3
    \end{cases}
    \]
    解得: $A_{-1} = \frac{h}{3}, A_1 = \frac{h}{3}, A_0 = \frac{4h}{3}$.
    此即辛普森公式, 代数精度为 3 (对 $x^3$ 也成立, 对 $x^4$ 不成立).

    \item[2.] 分别用梯形公式和辛普森公式计算 $\int_0^1 \frac{x}{4+x^2} \ud x, n=8$.
    
    \sol 准确值 $I = \frac{1}{2}[\ln(4+x^2)]_0^1 = \frac{1}{2}(\ln 5 - \ln 4) \approx 0.111572$.
    $h = 1/8 = 0.125$. 节点 $x_i = 0, 0.125, \dots, 1$.
    梯形公式: $T_8 = \frac{0.125}{2}[f(0) + 2\sum_{i=1}^7 f(x_i) + f(1)] \approx 0.1114$.
    辛普森公式: $S_8 \approx 0.11157$.

    \item[6.] 计算 $\int_0^1 \e^x \ud x$, 误差 $\leqslant 0.5 \times 10^{-5}$.
    
    \sol 
    (1) 梯形公式误差限 $|R_T| \leqslant \frac{h^2}{12}(b-a) M_2$.
    $M_2 = \max_{[0,1]} |\e^x| = \e$.
    $\frac{1}{12 n^2} \cdot \e \leqslant 0.5 \times 10^{-5} \implies n^2 \geqslant \frac{\e}{6 \times 10^{-5}} \approx 45304 \implies n \geqslant 213$.
    
    (2) 辛普森公式误差限 $|R_S| \leqslant \frac{h^4}{180}(b-a) M_4$.
    $M_4 = \e$.
    $\frac{1}{180 n^4} \cdot \e \leqslant 0.5 \times 10^{-5} \implies n^4 \geqslant \frac{\e}{90 \times 10^{-5}} \approx 3020 \implies n \geqslant 8$. (注意辛普森公式 $n$ 需为偶数分段或节点数配合).

    \item[8.] 龙贝格求积计算 $\frac{2}{\sqrt{\pi}} \int_0^1 \e^{-x} \ud x$.
    
    \sol
    $T_1 = \frac{1}{2}[f(0)+f(1)] \dots$
    $T_2 = \frac{1}{2}T_1 + \frac{1}{2}f(0.5) \dots$
    加速公式: $S_k = \frac{4T_{2k}-T_k}{3}$, $C_k = \frac{16S_{2k}-S_k}{15}$, $R_k = \frac{64C_{2k}-C_k}{63}$.
    最终结果 $I \approx 0.7132717$.

    \item[10.] 构造高斯型求积公式 $\int_0^1 \frac{1}{\sqrt{x}} f(x) \ud x \approx A_0 f(x_0) + A_1 f(x_1)$.
    
    \sol 权函数 $\rho(x) = x^{-1/2}$.
    矩 $\mu_k = \int_0^1 x^k x^{-1/2} \ud x = \int_0^1 x^{k-0.5} \ud x = \frac{1}{k+0.5}$.
    $\mu_0=2, \mu_1=2/3, \mu_2=2/5, \mu_3=2/7$.
    构造正交多项式 $x^2+ax+b$ 与 $1, x$ 正交:
    $\begin{cases} \mu_2 + a\mu_1 + b\mu_0 = 0 \\ \mu_3 + a\mu_2 + b\mu_1 = 0 \end{cases} \implies \begin{cases} 2/5 + 2a/3 + 2b = 0 \\ 2/7 + 2a/5 + 2b/3 = 0 \end{cases}$.
    解得 $a = -6/7, b = 3/35$.
    方程 $x^2 - \frac{6}{7}x + \frac{3}{35} = 0$ 的根为节点 $x_{0,1} = \frac{3 \mp \sqrt{9 - 8.4}}{7} = \frac{3 \mp \sqrt{0.6}}{7}$.
    系数通过代数精度方程求解 $A_0, A_1$.

    \item[14.] 计算 $\int_1^3 \frac{\ud y}{y}$. (3) 将积分区间分为四等份, 用复合两点高斯公式.
    
    \sol 子区间: $[1, 1.5], [1.5, 2], [2, 2.5], [2.5, 3]$.
    两点高斯公式在 $[a,b]$ 上的标准变换:
    $\int_a^b g(y)\ud y \approx \frac{b-a}{2} [g(y_1) + g(y_2)]$, 其中 $y_{1,2} = \frac{a+b}{2} \pm \frac{b-a}{2}\frac{1}{\sqrt{3}}$.
    
    (1) $[1, 1.5]$: 中点 1.25, 半长 0.25. 节点 $1.25 \pm \frac{0.25}{\sqrt{3}}$.
    $I_1 \approx 0.25 [\frac{1}{1.10566} + \frac{1}{1.39434}] \approx 0.405407$.
    
    (2) $[1.5, 2]$: 中点 1.75.
    $I_2 \approx 0.25 [\frac{1}{1.60566} + \frac{1}{1.89434}] \approx 0.287683$.
    
    (3) $[2, 2.5]$: 中点 2.25.
    $I_3 \approx 0.25 [\frac{1}{2.10566} + \frac{1}{2.39434}] \approx 0.223144$.
    
    (4) $[2.5, 3]$: 中点 2.75.
    $I_4 \approx 0.25 [\frac{1}{2.60566} + \frac{1}{2.89434}] \approx 0.182322$.
    
    总和 $I \approx 1.098556$. (精确值 $\ln 3 = 1.098612$).

    \item[18.] 用三点公式求导 $f(x) = (1+x)^{-2}$.
    \sol 
    $f'(x_0) \approx \frac{1}{2h}[-3f_0 + 4f_1 - f_2] \approx -0.247$.
    $f'(x_1) \approx \frac{1}{2h}[f_2 - f_0] \approx -0.217$.
    $f'(x_2) \approx \frac{1}{2h}[f_0 - 4f_1 + 3f_2] \approx -0.187$.
\end{enumerate}

% ==================== 第五章 ====================
\section*{第五章}

\begin{enumerate}[label=\textbf{5-\arabic*.}, leftmargin=*]
    \setcounter{enumi}{7}
    \item 用列主元消去法解线性方程组并求行列式.
    \[
    \bm{A} = \begin{bmatrix} 12 & -3 & 3 \\ -18 & 3 & -1 \\ 1 & 1 & 1 \end{bmatrix}, \bm{b} = \begin{bmatrix} 15 \\ -15 \\ 6 \end{bmatrix}.
    \]
    
    \sol
    1. 第一列主元选择: $|-18| > |12| > |1|$. 交换 R1, R2.
    \[
    \begin{bmatrix} -18 & 3 & -1 & -15 \\ 12 & -3 & 3 & 15 \\ 1 & 1 & 1 & 6 \end{bmatrix}.
    \]
    消元:
    $R_2 = R_2 - (12/-18)R_1 = R_2 + (2/3)R_1 \to [0, -1, 7/3, 5]$.
    $R_3 = R_3 - (1/-18)R_1 = R_3 + (1/18)R_1 \to [0, 7/6, 17/18, 31/6]$.
    
    2. 第二列主元选择: $|7/6| > |-1|$. 交换 R2, R3.
    \[
    \begin{bmatrix} -18 & 3 & -1 & -15 \\ 0 & 7/6 & 17/18 & 31/6 \\ 0 & -1 & 7/3 & 5 \end{bmatrix}.
    \]
    消元: $R_3 = R_3 - (-1 / (7/6))R_2 = R_3 + (6/7)R_2$.
    $a_{33} = 7/3 + (6/7)(17/18) = 7/3 + 17/21 = 22/7$.
    $b_3 = 5 + (6/7)(31/6) = 5 + 31/7 = 66/7$.
    
    3. 回代:
    $x_3 = (66/7) / (22/7) = 3$.
    $x_2 = (31/6 - (17/18)*3) / (7/6) = 2$.
    $x_1 = (-15 - 3*2 + 1*3) / (-18) = 1$.
    解为 $\bm{x} = [1, 2, 3]^\T$.
    
    行列式: $\det \bm{A} = (-1)^2 \times (-18) \times (7/6) \times (22/7) = -66$.

    \item[9.] 追赶法解三对角方程组.
    \sol 
    $L$ 对角元 $\alpha_i$, $U$ 上对角元 $\beta_i$.
    $\alpha_1 = 2, \beta_1 = -1/2$.
    $\alpha_2 = 2 - (-1)(-0.5) = 1.5, \beta_2 = -2/3$.
    $\alpha_3 = 4/3, \beta_3 = -3/4$.
    $\alpha_4 = 5/4, \beta_4 = -4/5$.
    $\alpha_5 = 6/5$.
    解 $Ly = b$: $y = [1/2, 1/3, 1/4, 1/5, 1/6]^\T$.
    解 $Ux = y \implies x_5 = 1/6, x_4 = 1/3, x_3 = 1/2, x_2 = 2/3, x_1 = 5/6$.

    \item[11.] 判断矩阵 LU 分解.
    \sol
    A: $\Delta_2 = 1\times 4 - 2\times 2 = 0$, 不能直接分解.
    B: $\Delta_2 = 0$, 不能直接分解.
    C: 各阶主子式均不为0, 可以唯一分解.

    \item[12.] 计算 $\bm{A} = \begin{pmatrix} 0.6 & 0.5 \\ 0.1 & 0.3 \end{pmatrix}$ 的范数.
    
    \sol
    $\|\bm{A}\|_\infty = \max(1.1, 0.4) = 1.1$.
    $\|\bm{A}\|_1 = \max(0.7, 0.8) = 0.8$.
    $\|\bm{A}\|_F = \sqrt{0.36+0.25+0.01+0.09} = \sqrt{0.71} \approx 0.8426$.
    $\bm{A}^\T \bm{A} = \begin{pmatrix} 0.37 & 0.33 \\ 0.33 & 0.34 \end{pmatrix}$. 特征值 $\approx 0.6853, 0.0247$.
    $\|\bm{A}\|_2 = \sqrt{\lambda_{\max}(\bm{A}^\T \bm{A})} \approx 0.8279$.

    \item[18.] 计算 $\bm{A} = \begin{pmatrix} 100 & 99 \\ 99 & 98 \end{pmatrix}$ 的条件数.
    \sol
    $\bm{A}^{-1} = \begin{pmatrix} -98 & 99 \\ 99 & -100 \end{pmatrix}$.
    $\|\bm{A}\|_\infty = 199, \|\bm{A}^{-1}\|_\infty = 199 \implies \text{cond}_\infty = 39601$.
    $\lambda_{\max}(\bm{A}^\T\bm{A}) \approx 39206, \lambda_{\min} \approx 2.55\times 10^{-5}$.
    $\text{cond}_2 = \sqrt{\lambda_{\max}/\lambda_{\min}} \approx 39206$.

    \item[19.] 证明正交矩阵 $\bm{Q}$ 的 2-条件数为 1.
    \proof $\|\bm{Q}\|_2 = \sqrt{\rho(\bm{Q}^\T\bm{Q})} = \sqrt{\rho(\bm{I})} = 1$.
    同理 $\|\bm{Q}^{-1}\|_2 = \|\bm{Q}^\T\|_2 = 1$.
    故 cond $= 1 \times 1 = 1$.
\end{enumerate}
% ==================== 第六章 ====================
\section*{第六章}

\begin{enumerate}[label=\textbf{6-\arabic*.}, leftmargin=*]
    \item 考察 Jacobi/GS 收敛性并求解.
    \sol 系数矩阵严格对角占优, 故均收敛.
    迭代求解过程略(纯数值代入).

    \item[4.] $\bm{A} = \begin{pmatrix} 10 & a & 0 \\ b & 10 & b \\ 0 & a & 5 \end{pmatrix}$, 求迭代收敛条件.
    
    \sol
    Jacobi 迭代矩阵 $B_J = D^{-1}(L+U)$.
    特征方程 $|\lambda I - B_J| = 0 \implies \lambda(\lambda^2 - \frac{ab}{100} - \frac{ab}{50}) = 0$.
    $\lambda^2 = \frac{3ab}{100}$. 收敛要求 $|\lambda| < 1 \implies \frac{3|ab|}{100} < 1 \implies |ab| < 100/3$.
    Gauss-Seidel 同理, 得到相同收敛条件.

    \item[6.] 证明 $\bm{A} = \begin{pmatrix} 3 & 0 & -2 \\ 0 & 2 & 1 \\ -2 & 1 & 2 \end{pmatrix}$ 两种方法均收敛.
    \sol 计算 $\rho(B_J) = \sqrt{11/12} < 1$ 及 $\rho(B_{GS}) = 11/12 < 1$.
    $\rho(B_{GS}) < \rho(B_J)$, 故 GS 收敛更快.
\end{enumerate}

% ==================== 第七章 ====================
\section*{第七章}

\begin{enumerate}[label=\textbf{7-\arabic*.}, leftmargin=*]
    \item 用二分法求 $x^2 - x - 1 = 0$ 的正根 (误差 $< 0.05$).
    
    \sol $f(1)=-1, f(2)=1$.
    \begin{center}
    \begin{tabular}{|c|c|c|c|c|}
    \hline
    $k$ & $a_k$ & $b_k$ & $x_k$ & $f(x_k)$符号 \\
    \hline
    0 & 1 & 2 & 1.5 & - \\
    1 & 1.5 & 2 & 1.75 & + \\
    2 & 1.5 & 1.75 & 1.625 & + \\
    3 & 1.5 & 1.625 & 1.5625 & - \\
    4 & 1.5625 & 1.625 & 1.59375 & - \\
    5 & 1.59375 & 1.625 & 1.609375 & - \\
    \hline
    \end{tabular}
    \end{center}
    $x_5$ 时区间半径 $(1.625-1.59375)/2 \approx 0.015 < 0.05$. 根 $\approx 1.609$.

    \item[6.] 确定参数使迭代法 3 阶收敛.
    \sol 由 $\varphi(x^*) = x^*, \varphi'(x^*) = 0, \varphi''(x^*) = 0$ 条件导出的:
    $p(x) = \frac{1}{f'(x)}, \quad q(x) = \frac{f''(x)}{2[f'(x)]^3}$.

    \item[13.] 导出求 $\sqrt[3]{a}$ 的迭代公式 (基于 $f(x)=1-a/x^3=0$), 并计算 $\sqrt[3]{115}$.
    
    \sol 
    方程 $1 - a/x^3 = 0$. 导数 $f'(x) = 3a/x^4$.
    牛顿迭代: $x_{k+1} = x_k - \frac{1-a/x_k^3}{3a/x_k^4} = x_k - \frac{x_k^4 - ax_k}{3a} = \frac{3ax_k - x_k^4 + ax_k}{3a} = \frac{4}{3}x_k - \frac{x_k^4}{3a}$.
    
    计算 $\sqrt[3]{115}$ ($a=115$). 真值 $\approx 4.86$.
    取 $x_0 = 5$ (注意: 若取10, 第一步会变成 $13.33 - 10000/345 < 0$, 发散).
    $k=0: x_1 = \frac{4}{3}(5) - \frac{625}{345} \approx 6.6667 - 1.8116 = 4.8551$.
    $k=1: x_2 = \frac{4}{3}(4.8551) - \frac{555.66}{345} \approx 6.4735 - 1.6106 = 4.8629$.
    $k=2: x_3 \approx 4.8629$. (已稳定).
\end{enumerate}

% ==================== 第九章 ====================
\section*{第九章}

\begin{enumerate}[label=\arabic*., leftmargin=*]
    \item 用欧拉法解初值问题
    \[
    y' = x^2 + 100y^2, \quad y(0) = 0.
    \]
    取步长 $h=0.1$, 计算到 $x=0.3$ (保留到小数点后 4 位).
    
    \sol 欧拉法公式为
    \[
    y_{n+1} = y_n + h f(x_n, y_n) = y_n + h(x_n^2 + 100y_n^2), \quad n=0, 1, 2.
    \]
    代入 $y_0 = 0$, 计算结果为
    \[
    y(0.1) \approx y_1 = 0, \quad y(0.2) \approx y_2 = 0.0010, \quad y(0.3) \approx y_3 = 0.0050.
    \]

    \item 用改进欧拉法和梯形法解初值问题
    \[
    y' = x^2 + x - y, \quad y(0) = 0.
    \]
    取步长 $h=0.1$, 计算到 $x=0.5$, 并与准确解 $y = -\e^{-x} + x^2 - x + 1$ 相比较.
    
    \sol 改进欧拉法为
    \[
    y_{n+1} = y_n + \frac{h}{2}[f(x_n, y_n) + f(x_{n+1}, y_n + hf(x_n, y_n))],
    \]
    将 $f(x, y) = x^2 + x - y$ 代入, 得
    \[
    y_{n+1} = \left( 1 - h + \frac{h^2}{2} \right)y_n + \frac{h}{2}[(1-h)x_n(1+x_n) + (1+x_{n+1})x_{n+1}].
    \]
    梯形法公式为
    \[
    y_{n+1} = y_n + \frac{h}{2}[f(x_n, y_n) + f(x_{n+1}, y_{n+1})],
    \]
    将 $f(x, y) = x^2 + x - y$ 代入, 得
    \[
    y_{n+1} = \frac{2-h}{2+h}y_n + \frac{h}{2+h}[x_n(1+x_n) + x_{n+1}(1+x_{n+1})],
    \]
    将 $y_0=0, h=0.1$ 代入计算公式, 结果如下:
    \begin{center}
    \begin{tabular}{c|c|c|c|c}
    \hline
    $x_n$ & 改进欧拉法 $y_n$ & $|y(x_n)-y_n|$ & 梯形法 $y_n$ & $|y(x_n)-y_n|$ \\
    \hline
    0.1 & 0.005 500 & $0.337\,418\,036 \times 10^{-3}$ & 0.005 238 095 & $0.755\,132\,781 \times 10^{-4}$ \\
    0.2 & 0.021 927 500 & $0.658\,253\,078 \times 10^{-3}$ & 0.021 405 896 & $0.136\,648\,778 \times 10^{-3}$ \\
    0.3 & 0.050 144 388 & $0.962\,608\,182 \times 10^{-3}$ & 0.049 367 239 & $0.185\,459\,653 \times 10^{-3}$ \\
    0.4 & 0.090 930 671 & $0.125\,071\,672 \times 10^{-2}$ & 0.089 903 692 & $0.223\,738\,443 \times 10^{-3}$ \\
    0.5 & 0.144 992 257 & $0.152\,291\,668 \times 10^{-2}$ & 0.143 722 388 & $0.253\,048\,087 \times 10^{-3}$ \\
    \hline
    \end{tabular}
    \end{center}
    可见梯形法比改进欧拉法精确.
\end{enumerate}

\end{document}